\documentclass[11pt,a4paper]{article}
\usepackage[utf8]{inputenc}
\usepackage[T1]{fontenc}
\usepackage[italian]{babel}
\usepackage{amsmath,amssymb,graphicx,booktabs,siunitx,hyperref}
\usepackage[margin=2.5cm]{geometry}
\graphicspath{{figures/}}

\title{ltburst: il tone-burst sagomato di Linkwitz\\\large Diario di studio e ricostruzione per la suite SEAM-LTM}
\author{Giuseppe Silvi --- SEAM}
\date{2026}

\begin{document}
\maketitle
\begin{abstract}
Questo documento \`e un diario di studio passo passo.
Ricostruisce il tone-burst sagomato di Siegfried Linkwitz a partire dalla letteratura originale, dalla teoria fino al plugin, e ne valida ogni stadio con grafici confrontabili con le figure dei lavori originali.
\end{abstract}
\tableofcontents

\section{Introduzione e motivazione}\label{sec:intro}

Il tone-burst sagomato è un segnale di test a banda stretta ottenuto moltiplicando una sinusoide continua per una finestra raised-cosine (Hanning), in modo che lo spettro risultante sia concentrato in una banda di circa un terzo di ottava attorno alla frequenza centrale~\cite{linkwitz1978narrowband, linkwitz1980shaped}.
Il numero di cicli della sinusoide rimane costante (tipicamente cinque) al variare della frequenza, cosicché la durata temporale della finestra si riduce con l'aumentare della frequenza e la larghezza di banda percentuale rimane costante: il segnale è pertanto a Q costante per definizione.
Questa proprietà conferisce al burst una risoluzione in frequenza pari a circa un terzo di ottava su tutto lo spettro udibile, analoga a quella dei filtri a banda costante utilizzati in psicoacustica.

Nella suite \textsc{seam-ltm}, \emph{ltburst} occupa una posizione particolare sotto tre aspetti distinti.
È il primo generatore di segnale della suite, che fino a questo momento è composta esclusivamente da processori e plugin di spazializzazione.
È il primo oggetto per cui la specifica Faust non preesiste, ma viene redatta come parte integrante di questo stesso lavoro di studio, e andrà a costituire la libreria \texttt{seam.linkwitz.lib} (prefisso \texttt{slw}), seguendo il modello già adottato per Gerzon, Moorer e Roads.
È infine il primo plugin della suite costruito come diario di studio sistematico a partire dalla letteratura primaria, con l'obiettivo di rendere ogni passaggio — dalla matematica al codice C++ — leggibile e verificabile da uno studente.

Un aspetto pedagogicamente rilevante riguarda il confronto tra la finestra di Hann continua e l'approssimazione a scalini quantizzati descritta da Linkwitz nel sistema hardware originale~\cite{linkwitz1980shaped}.
Nell'implementazione analogica, un attenuatore programmabile commuta il livello a ogni semiciclo della sinusoide, introducendo discontinuità di pendenza ai passaggi per lo zero.
Queste discontinuità generano armoniche dispari che limitano la dinamica del sistema di misura, un effetto assente nella finestra matematicamente esatta.
Il confronto tra le due varianti — la finestra ideale e l'approssimazione hardware — costituisce uno degli esercizi centrali di questo studio.

La motivazione più diretta per includere \emph{ltburst} nella suite è la sua utilità come strumento di validazione per le reti passa-tutto già sviluppate: il plugin \emph{hilbert} e il codificatore \emph{x2uhj}.
Linkwitz stesso osserva che il burst sagomato ``\emph{appears to be a very sensitive indicator of phase distortion in subjective evaluations of all-pass networks}''~\cite{linkwitz1980shaped}, confermando che questo tipo di segnale rivela distorsioni di fase che misure in frequenza a regime potrebbero mascherare.
Disporre di un generatore di burst calibrato e documentato permette quindi di sottoporre le reti all-pass della suite a una verifica sistematica nel dominio del tempo.

Infine, questo documento risponde a un'assenza concreta nella letteratura tecnica disponibile in lingua italiana: non esiste, a conoscenza dell'autore, alcun testo che accompagni un lettore italiano attraverso la teoria, la ricostruzione matematica e l'implementazione di un segnale di test di questo tipo a partire dalle fonti originali.

\section{Studio delle fonti}\label{sec:fonti}
Da scrivere nel Task 5.

\section{Teoria e matematica}\label{sec:teoria}
Da scrivere nel Task 6.

\section{Obiettivi di validazione}\label{sec:validazione}
Da scrivere nel Task 7.

\section{Ricostruzione in Faust}\label{sec:faust}
Sviluppata nella Fase 2.

\section{Porting in C++}\label{sec:cpp}
Sviluppata nella Fase 3.

\section{Il plugin}\label{sec:plugin}
Sviluppata nella Fase 4.

\section{Conclusioni e verso un paper}\label{sec:conclusioni}
Sviluppata a fine percorso.

\bibliographystyle{plain}
\bibliography{refs}
\end{document}
